\cvsection{Research Experience}
\entrytwo{Undergraduate Thesis}{2025}
\textit{[Department / Lab / University], [Supervisor Name]}\par
\begin{itemize}
  \item Conducted research on energy-efficient data transmission in Wireless Sensor Networks (WSNs), with emphasis on the unbalanced energy depletion caused by redundant transceiver transmissions of correlated environmental data, the DGS-WSN framework (integrating DBSCAN clustering, synchronized dual-side GRU prediction, and sleep-awake scheduling), and extending the network lifetime to over 20,000 rounds while achieving an 82.24 percent transmission avoidance ratio and a 100 percent Packet Delivery Ratio (PDR).
  \item Developed the DGS-WSN prediction and scheduling framework and evaluated it using network stability metrics (First Node Death, Half Node Death, Last Node Death), total energy consumption, average latency, and Transmission Avoidance Ratio (TAR) on a simulated 100×100 m homogeneous sensor field against LSTM, LMS-LSTM, and un-scheduled GRU benchmarks.
  \item Collaborated with [advisor/lab members/team] on experimental design, analysis, and interpretation of results.
  \item Prepared materials for a peer-reviewed conference manuscript published in the 2026 International Conference on Computing, Communication, Security and Intelligent Systems (IC3SIS), with focus on clarity, reproducibility, and scientific communication.
\end{itemize}

\vspace{0.15em}
\entrytwo{Final Year Thesis Project}{2024}
\textit{[Department / University / Supervisor]}\par
\begin{itemize}
  \item Investigated [problem statement] and identified the key technical and practical challenges in [specific domain].
  \item Designed and implemented a [model/system/framework] to address the problem using [technology stack or methodology].
  \item Evaluated performance through [experiments, datasets, validation strategy], with emphasis on accuracy, robustness, and efficiency.
  \item Documented the project, findings, and implementation decisions in a structured technical report.
\end{itemize}
